% !PRINTED 398
\weberpage{398}
Aus (7) ergiebt sich, dass, wenn $u$ wächst, auch $t$ wachsen
muss, dass also, wenn $u$ einen möglichst kleinen positiven Werth
hat, auch $t$ möglichst klein ist. Man kann also von einer kleinsten
positiven Lösung reden, die wir mit $T$, $U$ bezeichnen wollen. Man
erhält sie, wenn man in den Formeln (4), (5) $n$ möglichst klein,
also gleich der Gliederzahl $\nu$ der Periode setzt. Ist daher $\nu$
gerade, so ist nur die Gleichung (9), nicht die Gleichung (8)
lösbar; ist aber $\nu$ ungerade, so ist sowohl (8) als (9) lösbar. Da
dies nur von der Discriminante $D$, nicht von der besonderen
Zahl $\omega$ abhängen kann, so folgt,

\gesperrt{9. dass bei einer Discriminante die verschiedenen
Perioden entweder alle eine gerade oder alle
eine ungerade Gliederzahl enthalten.}

Ist $D$ gerade, also durch 4 theilbar, so muss auch $t$ gerade
sein, und die Gleichung (7) lässt sich Glied für Glied durch 4
theilen. Ist $D$ ungerade, so sind auch $t$, $u$ entweder beide gerade
oder beide ungerade; sind sie beide ungerade, so sind ihre Quadrate,
wie alle ungeraden Quadratzahlen, nach dem Modul 8 mit
1 congruent, also muss $D \equiv 5\ (\mathrm{mod}\ 8)$ sein. Aber es ist nicht
immer möglich, wenn $D \equiv 5\ (\mathrm{mod}\ 8)$ ist, die Gleichung (7)
durch ungerade Zahlen zu lösen\footnote{Vgl. \gesperrt{Cayley}, Note sur l'équation $x^2 - D y^2 = \pm 4$, $D \equiv 5\ (\mathrm{mod}\ 8)$
Crelle's Journal, Bd. 53. Mathematical papers, Vol. IV. Tafeln für die
Lösungen der Pell'schen Gleichung in der Form $y^2 = a x^2 + 1$ hat
\gesperrt{Degen} berechnet (Canon Pellianus, Havniae 1817). Auch in \gesperrt{Legendre}'s
Zahlentheorie (deutsch von \gesperrt{Maser}) findet sich eine solche Tafel.}. Ist $D \equiv 1\ (\mathrm{mod}\ 8)$, so
müssen $t$, $u$ gerade sein.

Die Beispiele des \S.\,126 geben, auf diese Weise behandelt,
die folgenden kleinsten positiven Lösungen der Pell'schen
Gleichung, wobei der Factor 4, wo es möglich ist, weggehoben ist.
\[
\begin{aligned}
5^2 - 29.1^2 &= -4\\
70^2 - 29.13^2 &= -1\\
170^2 - 19.39^2 &= 1\\
6^2 - 37.1^2 &= -1\\
35^2 - 34.6^2 &= 1.
\end{aligned}
\]
